% !TEX program = xelatex
\documentclass[11pt,a4paper]{article}
\usepackage{fontspec}
\setmainfont{Calibri}
\setmonofont[Scale=0.88]{Consolas}
\newfontfamily\symbolfont{Segoe UI Symbol}
\newfontfamily\cjkfont{Microsoft YaHei}[Scale=0.9]
\usepackage{polyglossia}\setdefaultlanguage{polish}
\usepackage[table]{xcolor}
\usepackage{booktabs,longtable,array,geometry,enumitem,graphicx,fancyhdr,caption}
\usepackage[most]{tcolorbox}
\PassOptionsToPackage{hyphens}{url}
\usepackage[hidelinks,unicode]{hyperref}
\emergencystretch=2.5em
\geometry{a4paper,margin=2.1cm,top=2.3cm,bottom=2.3cm}
\setlength{\parindent}{0pt}\setlength{\parskip}{5pt}
\definecolor{apteal}{HTML}{0E7C86}
\definecolor{apteallight}{HTML}{E6F4F5}
\definecolor{apgrey}{HTML}{F3F4F6}
\definecolor{apamber}{HTML}{B45309}
\definecolor{apamberlight}{HTML}{FEF3C7}
\definecolor{apred}{HTML}{B91C1C}
\definecolor{apredlight}{HTML}{FEE2E2}
\definecolor{apgreen}{HTML}{15803D}
\definecolor{apgreenlight}{HTML}{DCFCE7}
\renewcommand{\arraystretch}{1.25}
\captionsetup{font=small,labelfont=bf}
\pagestyle{fancy}\fancyhf{}
\fancyfoot[L]{\small AMPERE POINT — testy bez chmury · Q11 z modułem Shelly · v1 · 2026-09-28}
\fancyfoot[R]{\small\thepage}
\renewcommand{\headrulewidth}{0pt}
\newtcolorbox{uwaga}[1][Uwaga]{enhanced,breakable,colback=apamberlight,colframe=apamber,title={#1},fonttitle=\bfseries,boxrule=0.6pt,arc=2pt,left=6pt,right=6pt,top=4pt,bottom=4pt}
\newtcolorbox{info}[1][Jak to działa]{enhanced,breakable,colback=apteallight,colframe=apteal,title={#1},fonttitle=\bfseries,boxrule=0.6pt,arc=2pt,left=6pt,right=6pt,top=4pt,bottom=4pt}
\newtcolorbox{wazne}[1][Ważne]{enhanced,breakable,colback=apredlight,colframe=apred,title={#1},fonttitle=\bfseries,boxrule=0.6pt,arc=2pt,left=6pt,right=6pt,top=4pt,bottom=4pt}
\newtcolorbox{przyklad}[1][Przykład liczbowy]{enhanced,breakable,colback=apgreenlight,colframe=apgreen,title={#1},fonttitle=\bfseries,boxrule=0.6pt,arc=2pt,left=6pt,right=6pt,top=4pt,bottom=4pt}
\newtcolorbox{kodbox}{enhanced,breakable,colback=apgrey,colframe=apgrey,boxrule=0pt,arc=2pt,left=5pt,right=5pt,top=3pt,bottom=3pt,fontupper=\ttfamily\footnotesize}
\setlist{nosep,leftmargin=*}
\setcounter{secnumdepth}{2}
\begin{document}

{\LARGE\bfseries Testy bez chmury: Q11 z modułem Shelly\par}
\vspace{4pt}\textcolor{apteal}{\rule{\textwidth}{1.2pt}}
\emph{Poziom B: Wi-Fi i internet są, chmura Shelly wyłączona w module. Testy B1–B12 · AMPERE POINT · 28 września 2026 · oprac. Dawid Cekała}


\section*{W skrócie}

\begin{itemize}
\item \textbf{Poziom B} to obecny stan DevKitu: chmura wyłączona od 25.09, sieć biura z internetem. Nic nie trzeba przestawiać.
\item \textbf{Większość podstaw już sprawdziliśmy} (rozdział 1): sterowanie przez API, MQTT, OCPP przez most, godzina z internetu.
\item \textbf{Zostało:}
\begin{itemize}
\item czy da się sterować prosto, bez programowania: aplikacja Shelly, strona WWW i Home Assistant (B1–B3);
\item co moduł robi sam: harmonogram, webhook i pamięć (B4–B6);
\item jak działa API i gdzie są jego granice (B7–B12).
\end{itemize}
\item \textbf{Ryzyko:} małe. Jedynie hasło do API (B11) może odciąć nasze narzędzia, więc robimy je na końcu programu.
\item Poziomy łączności, pojęcia i porównanie z Tuya są w dokumencie 00, w tabeli 3.1.
\end{itemize}

\section{Co już sprawdziliśmy na tym poziomie}

\begin{longtable}{>{\raggedright\arraybackslash}p{3.93cm}>{\raggedright\arraybackslash}p{1.12cm}>{\raggedright\arraybackslash}p{1.65cm}>{\raggedright\arraybackslash}p{8.13cm}}
\toprule
\rowcolor{apteallight}\textbf{Funkcja} & \textbf{Wynik} & \textbf{Kiedy} & \textbf{Szczegóły} \\
\midrule\endhead
\bottomrule\endfoot
Limit prądu i włącznik przez HTTP i WebSocket & \leavevmode{\symbolfont\color{apgreen}\textbf{✔}} & 25.09 & sterownik potwierdza; najpierw odsyła starą wartość, po około 1 s nową \\
MQTT z własnym brokerem & \leavevmode{\symbolfont\color{apgreen}\textbf{✔}} & 25.09 & limit 10 A przez MQTT potwierdzony przez sterownik w tej samej sekundzie \\
OCPP przez most & \leavevmode{\symbolfont\color{apamber}\textbf{◐}} & 25.09 & rejestracja, stan, znak życia, pomiary na żądanie, limit 10 i 12 A, zdalny start; pełna transakcja czeka na obciążenie \\
Godzina z publicznego serwera czasu & \leavevmode{\symbolfont\color{apgreen}\textbf{✔}} & 25.09 & chmura wyłączona o 11:48; o 11:51 sterownik dostał poprawną godzinę \\
Ikona sieci na wyświetlaczu & \leavevmode{\symbolfont\color{apamber}\textbf{◐}} & 25.09 & moduł wysyła stan „4” (połączony z chmurą), choć chmura jest wyłączona \\
Dziennik pracy modułu przez sieć & \leavevmode{\symbolfont\color{apgreen}\textbf{✔}} & 24–25.09 & cała diagnostyka bez kabla USB \\
Aplikacja Shelly przy włączonej chmurze & \leavevmode{\symbolfont\color{apgreen}\textbf{✔}} & 25.09 & 10:52–10:54: włącznik 4 razy, limit 6 razy; każde polecenie przyszło przez chmurę Shelly (w logu „via SHC”). Po wyłączeniu chmury o 11:48 aplikacja nie wysłała ani jednego polecenia; wszystkie późniejsze pochodzą z naszych narzędzi na laptopie \\
Aktualizacja z portalu przez Bluetooth & \leavevmode{\symbolfont\color{apamber}\textbf{◐}} & 24.09 & działa; przesyłanie potrafi się przerwać; wgranie konfiguracji kasuje Wi-Fi modułu \\
\end{longtable}

\section{Przygotowanie}

\begin{itemize}
\item DevKit w sieci biura pod adresem 192.168.0.238, chmura wyłączona.
\item \textbf{Serwer czasu: publiczny} (\texttt{time.\allowbreak{}cloudflare.\allowbreak{}com}), tak jak u klienta z internetem. Od 25.09 moduł jest przestawiony na nasz lokalny serwer 192.168.0.57; przed testami B wracamy do publicznego. Lokalny serwer jest potrzebny dopiero w dokumencie 02 i w testach D5–D6.
\item Broker MQTT może działać; testy B go nie używają, a MQTT na tym poziomie sprawdziliśmy 25.09.
\item Monitor logu uruchomiony.
\item Symulator auta podłączony (B5, B7).
\item Na laptopie Python z bibliotekami do prostego serwera HTTP i WebSocket (B5, B8).
\end{itemize}

\section{Testy}

\subsection{B1. Aplikacja Shelly w tej samej sieci}

\textbf{Cel.} Sprawdzić, czy zwykły użytkownik steruje ładowarką z aplikacji, gdy chmura modułu jest wyłączona. 25.09 aplikacja sterowała wyłącznie przez chmurę. Źródło polecenia w logu rozstrzygnie, czy bez chmury przejdzie na sieć lokalną: będzie to adres telefonu w sieci biura.


\begin{longtable}{>{\raggedright\arraybackslash}p{3.00cm}>{\raggedright\arraybackslash}p{4.69cm}>{\raggedright\arraybackslash}p{7.59cm}}
\toprule
\rowcolor{apteallight}\textbf{Krok} & \textbf{Co robimy} & \textbf{Wynik oczekiwany} \\
\midrule\endhead
\bottomrule\endfoot
B1a & Telefon w sieci biura, aplikacja Shelly, otwieramy DevKit & zapisujemy, czy urządzenie jest dostępne, czy „offline” \\
B1b & Zmiana limitu prądu i włącznika z aplikacji & jeśli działa: w logu modułu widać źródło polecenia z sieci lokalnej, sterownik potwierdza \\
B1c & Ten sam telefon na danych komórkowych & oczekujemy, że nie działa, bo zdalny dostęp idzie przez chmurę; potwierdzamy \\
\end{longtable}

\subsection{B2. Strona WWW modułu na produkcie Q11}

\textbf{Cel.} Strona modułu to droga bez aplikacji i bez konta. Działała 24.09 na symulatorze; teraz sprawdzamy ją na produkcie Q11.


\begin{longtable}{>{\raggedright\arraybackslash}p{3.00cm}>{\raggedright\arraybackslash}p{4.73cm}>{\raggedright\arraybackslash}p{7.55cm}}
\toprule
\rowcolor{apteallight}\textbf{Krok} & \textbf{Co robimy} & \textbf{Wynik oczekiwany} \\
\midrule\endhead
\bottomrule\endfoot
B2a & Laptop: przeglądarka, adres \texttt{http:\allowbreak{}/\allowbreak{}/\allowbreak{}192.\allowbreak{}168.\allowbreak{}0.\allowbreak{}238} & strona się otwiera, widać pola ładowarki \\
B2b & Zmiana limitu i włącznika ze strony & sterownik potwierdza w 1 s \\
B2c & Telefon w sieci biura: to samo & działa na telefonie; zapisujemy, czy jest czytelna \\
B2d & Czy strona pyta o hasło & dziś nie pyta; wniosek do analizy GAP: każdy w sieci może sterować ładowarką, dopóki nie ma hasła (B11) \\
\end{longtable}

\subsection{B3. Home Assistant z oficjalną integracją Shelly}

\textbf{Cel.} Tuya w Home Assistant wymaga klucza lokalnego i naszego profilu. Sprawdzić, czy Shelly wchodzi do Home Assistant bez tego.


\begin{longtable}{>{\raggedright\arraybackslash}p{3.00cm}>{\raggedright\arraybackslash}p{4.42cm}>{\raggedright\arraybackslash}p{7.86cm}}
\toprule
\rowcolor{apteallight}\textbf{Krok} & \textbf{Co robimy} & \textbf{Wynik oczekiwany} \\
\midrule\endhead
\bottomrule\endfoot
B3a & W naszym Home Assistant: dodanie urządzenia Shelly pod adresem 192.168.0.238 & Home Assistant wykrywa moduł \\
B3b & Przegląd encji & zapisujemy, które pola ładowarki są widoczne: stan, limit, włącznik, fazy, energia i czas sesji \\
B3c & Zmiana limitu i włącznika z Home Assistant & sterownik potwierdza \\
B3d & Po teście & urządzenie zostaje albo je usuwamy, decyzja użytkownika. Dodanie zmienia konfigurację naszego Home Assistant \\
\end{longtable}

\subsection{B4. Harmonogram w module}

\textbf{Cel.} Harmonogram zapisany w module ma działać bez chmury i przetrwać restart.


\begin{longtable}{>{\raggedright\arraybackslash}p{3.00cm}>{\raggedright\arraybackslash}p{7.04cm}>{\raggedright\arraybackslash}p{5.24cm}}
\toprule
\rowcolor{apteallight}\textbf{Krok} & \textbf{Co robimy} & \textbf{Wynik oczekiwany} \\
\midrule\endhead
\bottomrule\endfoot
B4a & Dwa zadania: za 2 min limit 6 A, za 4 min 16 A & oba odpalają o czasie, sterownik potwierdza \\
B4b & Restart modułu poleceniem, bez zaniku zasilania; nowe zadanie za 2 min & zadania są po restarcie i działają \\
\end{longtable}

\subsection{B5. Webhook}

\textbf{Cel.} Moduł ma sam zawiadamiać nasz system o zdarzeniach. Przykład: „zaczęło się ładowanie” trafia do programu na laptopie.


\begin{longtable}{>{\raggedright\arraybackslash}p{3.00cm}>{\raggedright\arraybackslash}p{5.89cm}>{\raggedright\arraybackslash}p{6.39cm}}
\toprule
\rowcolor{apteallight}\textbf{Krok} & \textbf{Co robimy} & \textbf{Wynik oczekiwany} \\
\midrule\endhead
\bottomrule\endfoot
B5a & Lista zdarzeń, na które moduł umie wywołać webhook (polecenie, które nic nie zmienia) & wiemy, czy zmiana pól ładowarki (stan, limit, włącznik) jest zdarzeniem, czy są tylko zdarzenia systemowe \\
B5b & Nasłuch HTTP na laptopie; webhook na zmianę stanu ładowania; symulator: wolne {\symbolfont→} podłączone {\symbolfont→} ładuje & wywołanie przy każdej zmianie, z treścią zdarzenia. Jeśli B5a pokaże brak takiego zdarzenia: \leavevmode{\symbolfont\color{apred}\textbf{✘}} i powód \\
B5c & Restart modułu, powtórka B5b & webhook zostaje i działa \\
\end{longtable}

\subsection{B6. Pamięć klucz-wartość}

\textbf{Cel.} Pamięć, w której skrypt mógłby trzymać dane sesji, np. energię i czas, żeby nie ginęły przy restarcie.


\begin{longtable}{>{\raggedright\arraybackslash}p{3.00cm}>{\raggedright\arraybackslash}p{4.52cm}>{\raggedright\arraybackslash}p{7.76cm}}
\toprule
\rowcolor{apteallight}\textbf{Krok} & \textbf{Co robimy} & \textbf{Wynik oczekiwany} \\
\midrule\endhead
\bottomrule\endfoot
B6a & Zapis wartości testowej & odczyt daje to samo \\
B6b & Restart modułu & wartość zostaje; zanik zasilania sprawdza C8 \\
\end{longtable}

\subsection{B7. API: wartości spoza zakresu}

\textbf{Cel.} Ustalić, kto pilnuje granic: pole modułu, skrypt czy sterownik. Definicja Tuya dopuszcza limit 6–32 A, a Q11 ma 16 A. Z symulatorem bez obciążenia to bezpieczne.


\begin{longtable}{>{\raggedright\arraybackslash}p{3.00cm}>{\raggedright\arraybackslash}p{5.37cm}>{\raggedright\arraybackslash}p{6.91cm}}
\toprule
\rowcolor{apteallight}\textbf{Krok} & \textbf{Co robimy} & \textbf{Wynik oczekiwany / co zapisujemy} \\
\midrule\endhead
\bottomrule\endfoot
B7a & Limit 5 A, 17 A, 20 A, 32 A & odpowiedź API (błąd czy zgoda) i czy coś poszło do sterownika; co sterownik odesłał \\
B7b & Limit 10,5 A i tekst zamiast liczby & to samo \\
B7c & Nieznane pole i nieznane polecenie & kod błędu API \\
\end{longtable}

\subsection{B8. API: wychodzący WebSocket}

\textbf{Cel.} Moduł sam łączy się z naszym serwerem. Nie trzeba wtedy otwierać portów w sieci klienta, a to podstawa tłumacza OCPP na serwerze AmperePoint.


\begin{longtable}{>{\raggedright\arraybackslash}p{3.00cm}>{\raggedright\arraybackslash}p{4.90cm}>{\raggedright\arraybackslash}p{7.38cm}}
\toprule
\rowcolor{apteallight}\textbf{Krok} & \textbf{Co robimy} & \textbf{Wynik oczekiwany / co zapisujemy} \\
\midrule\endhead
\bottomrule\endfoot
B8a & Serwer testowy na laptopie; w module adres tego serwera & moduł łączy się sam; zapisujemy, co wysyła po połączeniu i przy zmianach \\
B8b & Serwer odsyła polecenie: limit 10 A & polecenie dociera, sterownik potwierdza \\
B8c & Wyłączenie serwera na 1 min & czy i po ilu sekundach moduł łączy się ponownie \\
\end{longtable}

\subsection{B9. API: połączenia naraz}

\textbf{Cel.} Most OCPP, monitor, strona WWW, Home Assistant i aplikacja mogą łączyć się jednocześnie. Trzeba znać granicę.


\begin{longtable}{>{\raggedright\arraybackslash}p{3.00cm}>{\raggedright\arraybackslash}p{5.52cm}>{\raggedright\arraybackslash}p{6.76cm}}
\toprule
\rowcolor{apteallight}\textbf{Krok} & \textbf{Co robimy} & \textbf{Wynik oczekiwany / co zapisujemy} \\
\midrule\endhead
\bottomrule\endfoot
B9a & Most, monitor i strona WWW otwarte; dokładamy kolejne połączenia WebSocket z laptopa & przy którym połączeniu moduł odmawia albo zwalnia; czy zrywa któreś z istniejących \\
\end{longtable}

\subsection{B10. API: kopia konfiguracji}

\begin{longtable}{>{\raggedright\arraybackslash}p{3.00cm}>{\raggedright\arraybackslash}p{4.59cm}>{\raggedright\arraybackslash}p{7.69cm}}
\toprule
\rowcolor{apteallight}\textbf{Krok} & \textbf{Co robimy} & \textbf{Wynik oczekiwany / co zapisujemy} \\
\midrule\endhead
\bottomrule\endfoot
B10a & Utworzenie i pobranie kopii konfiguracji na laptop & co zawiera. Jeśli hasło Wi-Fi, plik nie może trafić do repozytorium \\
\end{longtable}

\subsection{B11. API: hasło}

\textbf{Cel.} Zabezpieczyć moduł w sieci klienta i ustalić, które drogi hasło obejmuje.


\begin{longtable}{>{\raggedright\arraybackslash}p{3.00cm}>{\raggedright\arraybackslash}p{6.75cm}>{\raggedright\arraybackslash}p{5.53cm}}
\toprule
\rowcolor{apteallight}\textbf{Krok} & \textbf{Co robimy} & \textbf{Wynik oczekiwany / co zapisujemy} \\
\midrule\endhead
\bottomrule\endfoot
B11a & Ustawienie hasła & — \\
B11b & Próba bez hasła i z hasłem: HTTP, WebSocket, strona WWW, Bluetooth, MQTT & które drogi wymagają hasła \\
B11c & Zdjęcie hasła & narzędzia działają jak przedtem \\
\end{longtable}

\textbf{Ryzyko: średnie.} Most OCPP i monitor przestaną działać, dopóki nie dopiszemy logowania. Hasło zapisujemy w pliku projektu, poza repozytorium, bo zgubione oznacza reset fabryczny.


\subsection{B12. Kanał UDP, opcjonalnie}

\begin{longtable}{>{\raggedright\arraybackslash}p{3.00cm}>{\raggedright\arraybackslash}p{6.00cm}>{\raggedright\arraybackslash}p{6.28cm}}
\toprule
\rowcolor{apteallight}\textbf{Krok} & \textbf{Co robimy} & \textbf{Wynik oczekiwany / co zapisujemy} \\
\midrule\endhead
\bottomrule\endfoot
B12a & Włączenie kanału UDP, polecenie odczytu stanu z laptopa & czy działa i do czego mógłby służyć \\
\end{longtable}

\section{Kolejność i czas}

\begin{longtable}{>{\raggedright\arraybackslash}p{3.99cm}>{\raggedright\arraybackslash}p{7.31cm}>{\raggedright\arraybackslash}p{3.99cm}}
\toprule
\rowcolor{apteallight}\textbf{Kolejność} & \textbf{Testy} & \textbf{Czas} \\
\midrule\endhead
\bottomrule\endfoot
1 & B1–B3 & 50 min \\
2 & B4–B6 & 1 h \\
3 & B7–B10, B12 & 1 h 30 min \\
4 & B11, na końcu całego programu & 20 min \\
\end{longtable}

Razem około 3 h 40 min. Miejsce tych testów w kolejności całego programu: dokument 00, rozdział 5.


\section{Wyniki}

\begin{longtable}{>{\raggedright\arraybackslash}p{3.00cm}>{\raggedright\arraybackslash}p{2.62cm}>{\raggedright\arraybackslash}p{7.10cm}>{\raggedright\arraybackslash}p{2.11cm}}
\toprule
\rowcolor{apteallight}\textbf{Test} & \textbf{Wynik} & \textbf{Co zaobserwowano} & \textbf{Data} \\
\midrule\endhead
\bottomrule\endfoot
B1 aplikacja w tej samej sieci & \leavevmode{\symbolfont\color{apred}\textbf{✘}} wstępnie, do powtórzenia & B1a: telefon w sieci biura, aplikacja pokazuje DevKit jako „offline”. W logu jedno połączenie z 192.168.0.146 o 07:12:06, otwarte i zamknięte bez polecenia. To adres tego telefonu (potwierdzony w B2c), czyli aplikacja zajrzała do modułu lokalnie, ale nic nie wysłała. B1b niewykonalne, bo urządzenie jest „offline”. B1c wynika z B1a. Wniosek wstępny: bez chmury aplikacja Shelly nie steruje, drogą dla użytkownika zostaje strona WWW (B2). Do powtórzenia po wgraniu v4, przy wyłączonej chmurze, z ekranem urządzenia otwartym w aplikacji & 29.09 07:12 \\
B1 powtórka (v4.2/v4.3) & \leavevmode{\symbolfont\color{apred}\textbf{✘}} & Aplikacja dalej „offline”. Telefon 192.168.0.146 łączył się z modułem 4 razy (09:44, 11:53, 13:01:08, 13:01:20). Za każdym razem krótkie zapytanie i zamknięcie po 3–6 s, bez polecenia i bez stałego połączenia; w logu wygląda to tak samo jak zapytanie o dane urządzenia (\texttt{/shelly}) z laptopa. Sieć działa: HA, strona WWW i laptop sterują modułem. Wniosek: aplikacja znajduje moduł w sieci, ale steruje przez chmurę. Rozstrzygnie próba kontrolna: chmura włączona na 2–3 min; jeśli aplikacja ożyje, przyczyną jest brak chmury & 29.09 13:28 \\
B1 próba kontrolna & \leavevmode{\symbolfont\color{apgreen}\textbf{✔}} rozstrzygnięte & Chmura włączona 13:43–13:45 (bez restartu przy wyłączaniu). Aplikacja od razu „online”. Polecenia z aplikacji przyszły przez serwer Shelly (\texttt{SHC 34.\allowbreak{}40.\allowbreak{}19.\allowbreak{}158:\allowbreak{}6022}, nadawca \texttt{iotc\_\allowbreak{}291\_\allowbreak{}1-\allowbreak{}shelly-\allowbreak{}288-\allowbreak{}eu}), nie przez sieć lokalną: limit prądu 6 {\symbolfont→} 12 A (sterownik potwierdził), wyłączenie ładowania. \textbf{Wniosek: aplikacja Shelly steruje tylko przez chmurę; w sieci lokalnej jedynie sprawdza moduł.} Przy wyłączeniu ładowania sterownik odpowiedział dwa razy „wł.”, a pole pokazuje „wył.”; moduł uznał odpowiedź za „bez zmian” i nie przekazał jej skryptowi (ten sam problem co odrzucony tryb w B3c) & 29.09 13:45 \\
B2 strona WWW na Q11 & \leavevmode{\symbolfont\color{apgreen}\textbf{✔}} & B2a \leavevmode{\symbolfont\color{apgreen}\textbf{✔}}: strona pod 192.168.0.238 z laptopa otwiera się po około 5 s, tytuł „Ampere Point Q11 DevKit test”; widać włącznik, stan, fazy, moc, energię, czas sesji i limit z suwakiem 6–16 A. B2b \leavevmode{\symbolfont\color{apgreen}\textbf{✔}}: limit 12 {\symbolfont→} 10 A ze strony, 07:13:25 zapis do sterownika, 07:13:27 potwierdzenie (najpierw stara wartość 12). Włącznik ze strony: 07:13:37 wyłączony, sterownik potwierdził i przeszedł w stan „czeka”; 07:15:34 włączony, znów „ładuje”. B2d: strona nie pyta o hasło. Usterki: kafelki faz pokazują prąd „N/A A”, choć moduł ma w polu faz prąd 0 i napięcie L1 227 V; czas sesji 1072 min to błąd skryptu (liczy z zegara). B2c \leavevmode{\symbolfont\color{apgreen}\textbf{✔}}: telefon w sieci biura (192.168.0.146), strona czytelna. Z telefonu o 07:18:04 wyłączone ładowanie, o 07:18:12 limit 9 A, o 07:18:17 limit 14 A, o 07:18:19 włączone ładowanie; każde polecenie sterownik potwierdził w ciągu 1–2 s. Obserwacja: przy zmianie na 14 A moduł wysłał zapis do sterownika dwa razy w tej samej sekundzie, przy jednym poleceniu; nieszkodliwe, do wyjaśnienia & 29.09 07:13–07:18 \\
B3 Home Assistant & \leavevmode{\symbolfont\color{apamber}\textbf{◐}} wstrzymany & B3a \leavevmode{\symbolfont\color{apgreen}\textbf{✔}}: oficjalna integracja Shelly dodała „Q11 DevKit test” lokalnie; w logu połączenia przez WebSocket od biblioteki Home Assistant (\texttt{aios-…}), bez chmury. B3b \leavevmode{\symbolfont\color{apamber}\textbf{◐}}: widoczne encje: energia ×2, moc, moc, napięcie i prąd faz A–C (L1 227 V; prąd 0,00 A, czyli dane faz są poprawne, a „N/A” na stronie WWW to błąd strony). Encji sterowania (włącznik, limit, stan) jeszcze nie sprawdzaliśmy. Test wstrzymany do rozbudowy produktu o brakujące pola & 29.09 07:20 \\
B3b po v4.2 & \leavevmode{\symbolfont\color{apamber}\textbf{◐}} & 19 encji: 12 fazowych i 7 systemowych, żadnego z 14 pozostałych pól; „Wczytaj ponownie” nic nie zmienia. Przyczyna w kodzie integracji: pola usługi dostają encję tylko, gdy ich klucz jest na liście znanych kluczy. Fazy są na liście dla każdego urządzenia. Limit prądu i włącznik ładowania są na liście tylko dla ładowarki TopAC (kod produktu EVE01), a nasz kod to apq11dev. Stan, energia i czas sesji są na liście pod innymi nazwami (work\_state, energy\_charge, time\_charge) niż w szablonie Shelly (mode, session\_energy, session\_duration). Pozostałe nasze pola mają klucze nieznane integracji & 29.09 11:45 \\
B3 z łatką & \leavevmode{\symbolfont\color{apgreen}\textbf{✔}} & Kopia integracji Shelly w \texttt{custom\_\allowbreak{}components} z dopisanym Q11 (folder \texttt{DevKit\textbackslash{}ha\_\allowbreak{}integracja}, opis w \texttt{CZYTAJ.pdf}). Po restarcie 33 encje: dochodzi 14, w tym limit prądu, włącznik ładowania, tryb ładowania, limit energii, okno od/do, stan ładowarki, sygnał pojazdu, energia i czas sesji, temperatura, usterki. W logu brak błędów & 29.09 11:57 \\
B3c & \leavevmode{\symbolfont\color{apgreen}\textbf{✔}} z uwagami & Wszystkie polecenia z Home Assistant poszły lokalnie (WebSocket z 192.168.0.57) i sterownik je potwierdził: limit prądu 14 {\symbolfont→} 9 {\symbolfont→} 6 A, włącznik 6 razy, tryb, okno (1–1 odrzucone, 2–0 przyjęte), limit energii. Uwagi: (1) skrypt dwa razy wysłał starą wartość po nowej (limit energii 3 {\symbolfont→} 2, włącznik „wł.” bez polecenia), bo sterownik odpowiada z opóźnieniem i starą wartością; poprawka w skrypcie v4.3. (2) Odrzucenie trybu „do limitu energii” przy limicie 0 nie dotarło do pola: HA pokazywał tryb, którego sterownik nie przyjął. (3) Pole limitu energii w HA wysyła polecenie po każdym kliknięciu strzałki: 24 polecenia w 9 s & 29.09 12:04 \\
B3c po v4.3.1 & \leavevmode{\symbolfont\color{apgreen}\textbf{✔}} & Skrypt v4.3.1 wgrany 13:25 (usługa 13109-3815). Próba przez API modułu, tak jak wysyła HA: seria 8 {\symbolfont→} 6 w 0,3 s (powrót do wartości sterownika) — nic nie poszło; seria 7 {\symbolfont→} 9 {\symbolfont→} 7 w 0,6 s — jeden zapis 7, sterownik najpierw odesłał starą 6, potem 7, bez powtórnego zapisu; powrót do 6 — jeden zapis. Po starcie modułu brak komunikatów (v4.3 przy każdym starcie odrzucała jedną zmianę trybu) & 29.09 13:27 \\
B4 harmonogram w module & \leavevmode{\symbolfont\color{apgreen}\textbf{✔}} & B4a \leavevmode{\symbolfont\color{apgreen}\textbf{✔}}: dwa zadania (\texttt{Schedule.\allowbreak{}Create}, polecenie \texttt{Number.Set} limitu) odpaliły o 13:54:00 i 13:56:00 czasu modułu (czas z internetu, strefa Europe/Warsaw); skrypt wysłał po 1 s, sterownik potwierdził 6 A i 16 A. B4b \leavevmode{\symbolfont\color{apgreen}\textbf{✔}}: po restarcie poleceniem (13:56:45) oba zadania były na liście, nowe zadanie odpaliło o 13:59:00, sterownik potwierdził 12 A. Zmiana z harmonogramu przychodzi do skryptu bez pola źródła, skrypt traktuje ją jak zmianę z zewnątrz. Zadania testowe usunięte (codzienne) & 29.09 13:59 \\
B5 webhook & \leavevmode{\symbolfont\color{apgreen}\textbf{✔}} (B5b niżej) & B5a \leavevmode{\symbolfont\color{apgreen}\textbf{✔}}: \texttt{Webhook.\allowbreak{}ListSupported} ma zdarzenia pól: \texttt{number.\allowbreak{}change}, \texttt{enum.change}, \texttt{boolean.\allowbreak{}change}, \texttt{text.change}, \texttt{object.\allowbreak{}change}, więc zmiana stanu, limitu i włącznika może wywołać webhook. B5b: webhooki założone (stan z wartością w adresie, start ładowania z warunkiem \texttt{ev.\allowbreak{}value == "charger\_\allowbreak{}charging"}, limit, włącznik); odbiornik w kontenerze Docker na laptopie (port 8099; zapora Windows blokuje zwykłe programy). Na zmianę limitu wywołania dochodzą w ok. 1 s z wartością (\texttt{/limit?v=11}). Przejścia stanu z testera: do zrobienia. B5c \leavevmode{\symbolfont\color{apgreen}\textbf{✔}} dla limitu: webhooki są po restarcie i działają & 29.09 14:05 \\
B5b z testerem & \leavevmode{\symbolfont\color{apgreen}\textbf{✔}} & 30.09 12:41, z opóźnieniem ustawionym w menu ładowarki. Webhooki (czas laptopa): \texttt{/\allowbreak{}stan?v=charger\_\allowbreak{}charging} 12:41:12.8, \texttt{/\allowbreak{}ladowanie\_\allowbreak{}start} (z warunkiem) 12:41:13.9, \texttt{/\allowbreak{}stan?v=charger\_\allowbreak{}wait} 12:41:15.9, \texttt{/\allowbreak{}wlacznik?v=false} 12:41:38.8, \texttt{/\allowbreak{}stan?v=charger\_\allowbreak{}free} 12:41:46.1. Każda zmiana stanu dała wywołanie w około 1 s, warunek na start ładowania działa. Stanu „podłączone” nie było: sterownik przeszedł od razu do „ładuje”. Sygnał z auta przełączany 9 V {\symbolfont↔} 6 V (12:41:28–43) nie zmienił stanu z „czeka”, co pasuje do działającego opóźnienia z menu & 30.09 12:42 \\
B6 pamięć klucz-wartość & \leavevmode{\symbolfont\color{apgreen}\textbf{✔}} & B6a: \texttt{KVS.Set}/\texttt{KVS.Get} dają to samo. B6b: po restarcie poleceniem wartość zostaje. Zanik zasilania: C8 & 29.09 14:05 \\
B7 wartości spoza zakresu & \leavevmode{\symbolfont\color{apgreen}\textbf{✔}} & Granic pilnuje pole modułu: 5, 17, 20, 32 A {\symbolfont→} błąd -103 „out of bounds”, nic do sterownika. Tekst „abc” {\symbolfont→} -103; tekst „12” przyjęty jako liczba. 10,5 A przyjęte (moduł nie pilnuje kroku), skrypt wysłał 11 A, sterownik potwierdził. Nieznane pole {\symbolfont→} -105, nieznane polecenie {\symbolfont→} 404, opcja spoza listy {\symbolfont→} -103, zapis pola tylko do odczytu {\symbolfont→} -107 „Permission denied” & 29.09 14:06 \\
B8 wychodzący WebSocket & \leavevmode{\symbolfont\color{apgreen}\textbf{✔}} & Serwer testowy w kontenerze na laptopie (port 8098). B8a: moduł połączył się sam 25 s po restarcie, przysłał pełny stan, potem każdą zmianę pola. B8b: polecenie z serwera „limit 10 A” {\symbolfont→} odpowiedź w 150 ms, zapis do sterownika. B8c: po zniknięciu serwera ponowienia po 30 s, potem co 60 s; po powrocie serwera połączenie po 28 s. Połączenie wychodzące wyłączone po teście & 29.09 14:09 \\
B9 połączenia naraz & \leavevmode{\symbolfont\color{apgreen}\textbf{✔}} z uwagą & Poza HA, monitorem, stroną WWW i połączeniem wychodzącym moduł przyjął 6 kolejnych połączeń WebSocket; siódmego nie (brak odpowiedzi HTTP). Czasy odpowiedzi 110–170 ms. W chwili odmowy moduł zamknął połączenie Home Assistant, a HA połączył się ponownie: przy komplecie połączeń nowy klient może wyrzucić istniejącego & 29.09 14:10 \\
B10 kopia konfiguracji & \leavevmode{\symbolfont\color{apgreen}\textbf{✔}} z uwagą & \texttt{Sys.\allowbreak{}CreateBackup} + restart, potem \texttt{Sys.\allowbreak{}DownloadBackup} z parametrem \texttt{offset} (po 1024 B; bez niego wciąż ten sam kawałek). ZIP 9 kB, \textbf{niezaszyfrowany}: konfiguracja z hasłem Wi-Fi i tokenem chmury, pamięć klucz-wartość, webhooki, harmonogram, stan pól. Bez plików usługi (skrypt, ekran). Plik w \texttt{DevKit\textbackslash{}kopie\_\allowbreak{}konfiguracji} z ostrzeżeniem, poza repozytorium & 29.09 14:12 \\
B11 hasło & czeka na Dawida & Ustawienie hasła to wpisanie hasła; robi to Dawid &  \\
B13 MQTT z brokerem w internecie (dodatkowy) & \leavevmode{\symbolfont\color{apgreen}\textbf{✔}} & 01.10 11:43–11:48. Moduł z internetem (adres automatyczny), broker publiczny \texttt{test.\allowbreak{}mosquitto.\allowbreak{}org:\allowbreak{}1883}, losowy przedrostek tematów. Moduł połączył się z brokerem i utrzymał połączenie; jego wiadomości (online, stan, zdarzenia) dochodziły do laptopa przez internet. Polecenie limitu 10 A z laptopa przez broker: odpowiedź modułu po około 1,1 s, sterownik potwierdził; powrót do 12 A tak samo. Broker publiczny zrywał co drugie połączenie laptopa przy łączeniu (klient bez identyfikatora odrzucany, potem losowo); połączenie modułu stabilne. Po teście broker z powrotem na laptopie. Uwaga: w tym czasie chmura modułu była włączona (włączyło ją ponowne dodanie w aplikacji 30.09) & 01.10 11:48 \\
B12 kanał UDP & \leavevmode{\symbolfont\color{apgreen}\textbf{✔}} & Port 1010: odczyt 40–60 ms, zapis 135 ms, polecenie dochodzi do sterownika. Tylko pytanie–odpowiedź, bez zdarzeń. Bez hasła, więc po teście wyłączony & 29.09 14:13 \\
\end{longtable}

\end{document}